\section{Algoritmos Criptográficos Empleados en HTTPS}

El funcionamiento seguro del protocolo HTTPS depende de un conjunto de algoritmos criptográficos que trabajan de forma coordinada para garantizar confidencialidad, integridad, autenticación y resistencia ante ataques avanzados. Estos algoritmos, seleccionados y estandarizados por entidades como el IETF y el NIST, son implementados durante el \textit{handshake} y en la protección de los datos transportados mediante \textit{Transport Layer Security} (TLS). Su correcta elección e implementación se consideran factores críticos para la solidez del canal seguro \cite{rfc8446, rfc7525}.

En esta sección se describen los algoritmos fundamentales utilizados en HTTPS, incluyendo sus propiedades, su rol dentro del protocolo y su relevancia en el ecosistema criptográfico contemporáneo.

\subsection{Cifrado Simétrico}

El cifrado simétrico constituye el mecanismo principal para garantizar la \textbf{confidencialidad} de los datos transmitidos una vez establecida la sesión HTTPS. Entre los algoritmos recomendados y empleados en TLS 1.3 destacan:

\begin{itemize}
    \item \textbf{AES-GCM (Advanced Encryption Standard - Galois/Counter Mode)}: Proporciona cifrado autenticado y es la opción preferida en TLS debido a su eficiencia y resistencia frente a ataques criptográficos avanzados \cite{rfc8446}. Al combinar cifrado en modo contador y autenticación mediante Galois, reduce la necesidad de mecanismos adicionales para garantizar integridad.
    
    \item \textbf{ChaCha20-Poly1305}: Algoritmo de cifrado autenticado diseñado para ser eficiente en dispositivos sin aceleración AES, como teléfonos móviles o sistemas embebidos \cite{krawczyk2016}. Su seguridad está sustentada en análisis matemáticos rigurosos y su rendimiento es consistente en arquitecturas heterogéneas.
\end{itemize}

TLS 1.3 elimina algoritmos considerados débiles o inseguros, como RC4 o 3DES, mitigando vulnerabilidades ampliamente documentadas en versiones previas del protocolo \cite{rfc7525}.

\subsection{Intercambio de Claves}

El intercambio de claves tiene como objetivo permitir que cliente y servidor acuerden material criptográfico sin exponerlo a un posible adversario que controle el canal. En TLS 1.3 este proceso se basa exclusivamente en mecanismos con \textbf{Perfect Forward Secrecy (PFS)}, lo cual garantiza que la seguridad de sesiones pasadas no se vea comprometida incluso si un atacante obtiene claves privadas del servidor.

Los algoritmos principales son:

\begin{itemize}
    \item \textbf{ECDHE (Elliptic Curve Diffie–Hellman Ephemeral)}: Esquema de intercambio de claves ampliamente adoptado debido a su eficiencia y a la fortaleza matemática de las curvas elípticas recomendadas por estándares como NIST P-256 o Curve25519 \cite{krawczyk2016}. Su naturaleza efímera proporciona PFS de forma nativa.
    
    \item \textbf{DHE (Diffie–Hellman Ephemeral)}: Variante tradicional basada en aritmética modular. Aunque más lenta que ECDHE, continúa siendo soportada por compatibilidad y también ofrece PFS.
\end{itemize}

Los estudios formales han validado la seguridad de estos esquemas bajo modelos criptográficos rigurosos, incluyendo escenarios adversariales avanzados \cite{dowling2015}.

\subsection{Algoritmos de Hash}

Las funciones hash se emplean en múltiples partes del protocolo TLS: generación de claves, autenticación de mensajes, verificación de integridad y construcción de estructuras criptográficas internas. Las funciones aprobadas por los estándares modernos incluyen:

\begin{itemize}
    \item \textbf{SHA-256} y \textbf{SHA-384}: Pertenecientes a la familia SHA-2, constituyen la base para la mayoría de las operaciones de autenticación y derivación de claves en TLS 1.3 \cite{rfc8446}. Su diseño resistente a colisiones, preimágenes y segundos preimágenes los posiciona como elementos centrales del ecosistema criptográfico actual.
\end{itemize}

TLS prohíbe el uso de funciones hash inseguras como SHA-1 debido a ataques de colisión demostrados experimentalmente.

\subsection{Firmas Digitales}

Las firmas digitales son utilizadas para autenticar la identidad del servidor, validar certificados X.509 y garantizar la integridad del handshake. Los algoritmos predominantes incluyen:

\begin{itemize}
    \item \textbf{RSA (Rivest–Shamir–Adleman)}: Tradicionalmente el algoritmo más utilizado para certificados TLS. Su seguridad se basa en la dificultad del problema de factorización de enteros. Para garantizar niveles de seguridad adecuados se requieren claves de al menos 2048 bits \cite{rfc5280}.
    
    \item \textbf{ECDSA (Elliptic Curve Digital Signature Algorithm)}: Alternativa moderna y eficiente basada en criptografía de curvas elípticas. Ofrece firmas cortas, validación rápida y alto nivel de seguridad con claves de menor tamaño en comparación con RSA \cite{krawczyk2016}.
\end{itemize}

Las recomendaciones contemporáneas favorecen el empleo de ECDSA en nuevos despliegues debido a su mejor desempeño y menor huella computacional.

\subsection{Mecanismos de Autenticación y MACs}

Para la autenticación de mensajes TLS emplea MACs (Message Authentication Codes) integrados dentro de los algoritmos AEAD (Authenticated Encryption with Associated Data). Los principales son:

\begin{itemize}
    \item \textbf{GCM Tag} (en AES-GCM): Resultado del proceso autenticado que valida integridad y autenticidad del mensaje.
    \item \textbf{Poly1305} (en ChaCha20-Poly1305): MAC altamente eficiente y matemáticamente robusto asociado al cifrado ChaCha20 \cite{krawczyk2016}.
\end{itemize}

El uso de AEAD simplifica la arquitectura del protocolo y elimina clases enteras de ataques asociados a MACs mal integrados.

\subsection{Síntesis de los Algoritmos Criptográficos}

En conjunto, los algoritmos criptográficos de HTTPS conforman un sistema robusto diseñado para enfrentar ataques contemporáneos y garantizar seguridad a largo plazo. TLS 1.3 adopta únicamente algoritmos modernos y elimina mecanismos vulnerables, lo que representa un avance significativo en la protección de comunicaciones digitales. Su diseño modular permite incorporar futuras mejoras criptográficas sin comprometer la interoperabilidad y la solidez matemática del ecosistema.

